\documentclass[16pt,openany]{book}

\usepackage[a4paper, margin=3cm]{geometry}
\usepackage{pagenote}
\usepackage{anyfontsize}
\usepackage[portuguese]{babel}
\usepackage{amsmath}    
\usepackage{amssymb}    
\usepackage{amsfonts}   
\usepackage{amsthm}
\usepackage{mathtools} 
\usepackage{bm}         
\usepackage{mathrsfs}   
\usepackage{esint}      
\usepackage{cancel}     
\usepackage{nicefrac}   
\usepackage{siunitx}    
\usepackage{array}      
\usepackage{arydshln}  
\usepackage{ytableau}   
\usepackage{tcolorbox}
\tcbuselibrary{breakable}
\usepackage{xcolor}
\usepackage{hyperref}
\usepackage{pict2e}
\usepackage{tikz}
\usepackage{tikz-cd}
\definecolor{lightred}{rgb}{1,0.6,0.6}
\usepackage{stmaryrd}
\usepackage{comment}
\usepackage{multicol}
\usepackage{enumitem}

\usepackage{tgpagella}
\usepackage{eulervm}
%\usepackage[euler-digits]{eulervm}

\usetikzlibrary{babel}
\theoremstyle{definition}

\newtheorem{definition}{Definição}[section]  
\newtheorem{example}{Exemplo}[section]       
\newtheorem{theorem}{Teorema}[section]     
\newtheorem{proposition}{Proposição}[section] 
\newtheorem{corollary}{Colorário}[section] 
\newtheorem{lemma}{Lema}[section]  


\newtcolorbox{mybox}[2]{
  colback=#1!10!white, 
  colframe=#1!50!black, 
  boxrule=0.5pt,        
  sharp corners,        
  fonttitle=\bfseries,  
  breakable,            
  title={\textbf{#2}},  
}


\renewenvironment{definition}[1][]{
  \refstepcounter{definition}
  \begin{mybox}{green}{Definição \thesection.\arabic{definition}}
}{\end{mybox}}

\renewenvironment{example}[1][]{
  \refstepcounter{example}
  \begin{mybox}{yellow}{Exemplo \thesection.\arabic{example}}
}{\end{mybox}}

\renewenvironment{theorem}[1][]{
  \refstepcounter{theorem}
  \begin{mybox}{blue}{Teorema \thesection.\arabic{theorem}}
}{\end{mybox}}

\renewenvironment{proposition}[1][]{
  \refstepcounter{proposition}
  \begin{mybox}{red}{Proposição \thesection.\arabic{proposition}}
}{\end{mybox}}

\renewenvironment{corollary}[1][]{
  \refstepcounter{corollary}
  \begin{mybox}{lightred}{Corolário \thesection.\arabic{corollary}}
}{\end{mybox}}

\renewenvironment{lemma}[1][]{
  \refstepcounter{lemma}
  \begin{mybox}{cyan}{Lema \thesection.\arabic{lemma}}
}{\end{mybox}}

\begin{comment}
\newcommand{\titlepagecontent}{
    \begin{center}
        \vspace*{2cm} 
        {\Large Apostila Álgebra Linear} \\[2cm]
        
        {\Large Marco Vieira Busetti} \\[2.5cm]
        
        {\Large Universidade Tecnológica Federal do Paraná} \\[0.5cm]
        {\Large Curitiba, Outubro de 2026}
    \end{center}
}
\end{comment}

\newcommand{\titlepagecontent}{
    \begin{tikzpicture}[remember picture, overlay]
        \node[opacity=0.5, inner sep=0pt, anchor=south east] 
        at ([xshift=6cm, yshift=-1cm]current page.south east) 
        {\includegraphics[width=1.4\textwidth]{imgs/peticov.jpg}};
    \end{tikzpicture}
    
    
    \vspace*{2cm}
    \begin{flushleft}
        \hspace*{0.5cm} 
        \begin{minipage}{0.8\textwidth} 
            {\Huge\bfseries Álgebra Linear} \\[1.2cm] 
            
            {\Large Marco Vieira Busetti} \\[4.2cm]
            
            %{\large Universidade Tecnológica Federal do Paraná} \\[0.5cm]
            %{\large Curitiba, Novembro de 2024}
        \end{minipage}
    \end{flushleft}
    \clearpage
}


\newcommand{\trianglechar}{
  \mathrel{
    \unitlength=1mm
    \begin{picture}(5,3)(0,-1.5)
      \polyline(0,0)(4,1.5)(4,-1.5)(0,0)
      \polyline(1.25,0)(3,0.75)(3,-0.75)(1.25,0)
    \end{picture}
  }
}

\newcommand{\trianglechareq}{
  \mathrel{
    \unitlength=1mm
    \begin{picture}(5,3)(0,-1.5)
      \polyline(0,0)(4,1.5)(4,-1.5)(0,0)
      \polyline(1.25,0)(3,0.75)(3,-0.75)(1.25,0)
      \polyline(0,-1.9)(4,-1.9)
    \end{picture}
  }
}

\newtheoremstyle{no-italic}  
  {3pt}   
  {3pt}   
  {\normalfont} 
  {}      
  {\bfseries} 
  {.}     
  { }     
  {}


\begin{document}

\begin{titlepage}
    \titlepagecontent
\end{titlepage}

\fontsize{12}{14}\selectfont

\tableofcontents


\chapter*{Lista de Símbolos e Notações}
\addcontentsline{toc}{chapter}{Lista de Símbolos e Notações}
\markboth{LISTA DE SÍMBOLOS E NOTAÇÕES}{LISTA DE SÍMBOLOS E NOTAÇÕES}

Abaixo estão listadas as principais notações utilizadas ao longo deste trabalho.
\vspace{0.8cm}

\begin{comment}
\setlength{\columnsep}{1.5cm} 
\begin{multicols}{2}

%    \begin{center}\textbf{Geral e Conjuntos}\end{center}
    \begin{description}[leftmargin=2.2cm, labelwidth=1.8cm, labelsep=0.4cm, itemsep=0.3cm, font=\normalfont]
        \item[{ $\mathbb{N}$ }] Conjunto dos números naturais: $\{0, 1, 2,\ldots\}$
        \item[{ $\mathbb{Z}$ }] Conjunto dos números inteiros: $\{\ldots, -1,0,1,2,\ldots\}$
        \item[{ $\mathbb{Q}$ }] Conjunto dos números racionais
        \item[{ $\mathbb{R}$ }] Conjunto dos números reais
        \item[{ $\mathbb{C}$ }] Conjunto dos números complexos
        \item[{ $\mathbb{U}$ }] Círculo unitário (conjunto dos números complexos de módulo 1)
        \item[{ $X^*$ }] Conjunto $X$ excluindo o elemento nulo \\(ex: $\mathbb{R}^* = \mathbb{R} \setminus \{0\}$)
        \item[{ $X_+$ }] Subconjunto dos elementos não negativos de $X$ (ex: $\mathbb{Z}_+ = \mathbb{N}$)
        \item[{ $X_{+}^{*}$ }] Subconjunto dos elementos estritamente positivos de $X$
        \item[{ $X_{\geq n}$ }] Subconjunto dos elementos de $X$ que são maiores ou iguais a $n$ (ex: $\mathbb{Z}_{\geq 2} = \{2,3,4,\ldots\}$)
        \item[{ $X_{< n}$ }] Subconjunto dos elementos de $X$ que são estritamente menores que $n$
        \item[{ $A \times B$ }] Produto cartesiano entre os conjuntos $A$ e $B$
        \item[{ $\text{Card}(X)$ }] Cardinalidade do conjunto $X$
        \item[{ $\llbracket a, b\rrbracket$ }] Conjunto de inteiros compreendidos entre $a$ e $b$
        \item[{ $a \mid b$ }] O inteiro $a$ divide o inteiro $b$
        \item[{ $a \nmid b$ }] O inteiro $a$ não divide o inteiro $b$
    \end{description}
    
    %\columnbreak
    
    \begin{description}[leftmargin=2.2cm, labelwidth=1.8cm, labelsep=0.4cm, itemsep=0.3cm, font=\normalfont]
        \item[{ $e$ ou $e_G$}] Elemento neutro do grupo $G$
        \item[{ $|G|$ }] Ordem do grupo $G$
        \item[{ $H \le G$ }] $H$ é subgrupo de $G$
        \item[{ $H < G$ }] $H$ é subgrupo próprio de $G$
        \item[{ $N \trianglelefteq G$ }] $N$ é subgrupo normal de $G$
        \item[{ $N \triangleleft G$ }] $N$ é subgrupo normal próprio de $G$
        \item[{ $H \trianglechareq G$ }] $H$ é subgrupo característico de $G$
        \item[{ $[G:H]$ }] Índice do subgrupo $H$ em $G$
        \item[{ $G/H$ }] Grupo quociente de $G$ por $H$
        \item[{ $Z(G)$ }] Centro do grupo $G$
        \item[{ $\langle X \rangle$ }] Subgrupo gerado pelo conjunto $X$
        \item[{ $[x,y]$ }] Comutador dos elementos $x$ e $y$
        \item[{ $G'$ ou $G^{(1)}$ }] Grupo derivado (ou dos comutadores) de $G$
        \item[{ $G^{ab}$ }] Abelianização de $G$ (quociente $G/G'$)
        \item[{ $\mathcal{O}(g)$ }] Ordem do elemento $g$
        \item[{ $\ker \varphi$ }] Núcleo do homomorfismo $\varphi$
        \item[{ $\operatorname{Im} \varphi$ }] Imagem do homomorfismo $\varphi$
        \item[{ $\operatorname{Hom}(G,\mathcal{G})$ }] Conjunto de homomorfismos de $G$ em $\mathcal{G}$
        \item[{ $\operatorname{Aut}(G)$ }] Grupo de automorfismos de $G$
        \item[{ $\operatorname{Inn}(G)$ }] Grupo de automorfismos internos de $G$
        \item[{ $G \cong \mathcal{G}$ }] Isomorfismo entre $G$ e $\mathcal{G}$
    \end{description}

%    \begin{center}\textbf{Produtos de Grupos}\end{center}
    \begin{description}[leftmargin=2.2cm, labelwidth=1.8cm, labelsep=0.4cm, itemsep=0.3cm, font=\normalfont]
        \item[{ $\prod G_i$ }] Produto direto externo de grupos
        \item[{ $\bigodot H_i$ }] Produto direto interno de grupos
        \item[{ $H \rtimes K$ }] Produto semidireto interno
        \item[{ $H \rtimes_\theta K$ }] Produto semidireto externo associado a $\theta$
    \end{description}
    
    %\columnbreak

%    \begin{center}\textbf{Grupos Específicos}\end{center}
    \begin{description}[leftmargin=2.2cm, labelwidth=1.8cm, labelsep=0.4cm, itemsep=0.3cm, font=\normalfont]
        \item[{ $GL_n(\mathbb{K})$ }] Grupo linear geral sobre $\mathbb{K}$
        \item[{ $SL_n(\mathbb{K})$ }] Grupo linear especial sobre $\mathbb{K}$
        \item[{ $\mathbb{Z}/n\mathbb{Z}$ }] Grupo dos inteiros módulo $n$
        \item[{ $\mathbb{Z}_{p}^{*}$ }] Grupo multiplicativo dos inteiros invertíveis módulo $p$
        \item[{ $\mathfrak{S}_X$ ou $\mathfrak{S}_n$ }] Grupo simétrico do conjunto $X$ (ou de grau $n$)
        \item[{ $\mathfrak{A}_X$ ou $\mathfrak{A}_n$ }] Grupo alternado do conjunto $X$ (ou de grau $n$)
        \item[{ $\mathcal{K}$ }] Grupo de Klein
        \item[{ $\mathfrak{K}$ }] Grupo simétrico de Klein
    \end{description}
    
    %\begin{center}\textbf{Permutações e Ações}\end{center}
    \begin{description}[leftmargin=2.2cm, labelwidth=1.8cm, labelsep=0.4cm, itemsep=0.3cm, font=\normalfont]
        \item[{ $\text{id}_X$ }] Permutação identidade (ou função identidade)
        \item[{ $\text{supp}(\sigma)$ }] Suporte da permutação $\sigma$
        \item[{ $\mathscr{I}(\sigma)$ }] Número de inversões da permutação $\sigma$
        \item[{ $\varepsilon(\sigma)$ }] Assinatura da permutação $\sigma$
        \item[{ $G \curvearrowright X$ }] $G$ age à esquerda sobre $X$
        \item[{ $\mathscr{O}(x)$ }] Órbita do elemento $x$
        \item[{ $X/G$ }] Conjunto das órbitas da ação de $G$ sobre $X$
        \item[{ $G_x$ }] Estabilizador (grupo de isotropia) de $x$
        %\item[{ $\text{Fix}(X)$ }] Conjunto de pontos fixos da ação
    \end{description}

\end{multicols}

\end{comment}

\newpage

\chapter{Noções sobre Grupos, Anéis e Corpos}

\section{Definição e exemplos de grupos}

\section{Definição e exemplos de anéis}

\chapter{Generalidades sobre Espaços Vetoriais}

\section{Espaços Vetoriais}

\begin{definition}

Seja $(\mathbb{K}, \oplus, \odot)$ um corpo e $E$ um conjunto não vazio. Dizemos que $E$ munido das leis interna e externa, respectivamente:

\[
+\;:
\begin{array}{r c l}
  E\times E &\to &  E \\
  (x,y) &\mapsto & x+y
\end{array}
\qquad \text{e,} \qquad
\cdot\;:
\begin{array}{r c l}
  \mathbb{K}\times E &\to & E \\
  (\lambda, x) &\mapsto & \lambda x
\end{array}
\]

\textbf{é um espaço vetorial sobre $\mathbb{K}$} se, e somente se, $(E, +)$ é um grupo abeliano e:

\begin{itemize}
  \item[(i)] $\lambda(\mu x) = (\lambda \odot \mu)x, \quad \forall x\in E,\; \forall \lambda,\mu\in \mathbb{K}$
  \item[(ii)] $(\lambda \oplus \mu)x = \lambda x + \mu x, \quad \forall x\in E,\; \forall \lambda,\mu\in \mathbb{K}$
  \item[(iii)] $\lambda(x+y) = \lambda x + \lambda y, \quad \forall x,y\in E,\; \forall \lambda\in \mathbb{K}$
  \item[(iv)] $1\cdot x = x, \quad \forall x\in E$
\end{itemize}

onde $1$ representa o elemento neutro para a multiplicação em $\mathbb{K}$.
\end{definition}

\textit{Observações:}

\begin{itemize}
  \item \textit{Chamamos comumente os elementos de $\mathbb{K}$ de \textbf{escalares} e os elementos de $E$ de \textbf{vetores}};
  \item \textit{Denominamos um espaço vetorial $E$ sobre $\mathbb{K}$ de $\mathbb{K}$-espaço vetorial};
  \item \textit{Caso não haja ambiguidade, denotamos comumente os mesmos símbolos usuais de soma e produto para as operações do corpo $\mathbb{K}$ e as leis interna e externa do espaço vetorial $E$};
  \item \textit{Note que a definição de espaço vetorial é equivalente a afirmar que o corpo $\mathbb{K}$ \textbf{age} sobre o grupo abeliano $E$}.
\end{itemize}

\begin{proposition}

Sejam $\mathbb{K}$ um corpo e $E$ um $\mathbb{K}$-espaço vetorial. Então, $\forall (\lambda,x)\in \mathbb{K}\times E$, temos:

\begin{itemize}
  \item[(i)] $\lambda 0_E$ e $0x = 0_E$
  \item[(ii)] $\lambda x = 0_E \implies \lambda = 0$ ou $x = 0_E$
  \item[(iii)] $(-\lambda)x = \lambda(-x) = -(\lambda x)$
\end{itemize}

\end{proposition}

\textbf{Demonstração:}

Mostremos primeiramente o item $(i)$. Note que:

\[
\lambda 0_E = \lambda (0_E + 0_E) = \lambda 0_E + \lambda 0_E \implies \lambda 0_E = 0_E
\]

Analogamente, temos também que:

\[
0x = (0+0)x = 0x + 0x \implies 0x = 0_E 
\]

Note que o item $(ii)$ é trivial pela contrapositiva partindo do item $(i)$.

Mostremos, por fim, o item $(iii)$. Note que:

\[
(-\lambda + \lambda)x = 0x \iff (-\lambda)x + \lambda x = 0_E \iff (-\lambda)x = -(\lambda x)
\]

A prova é análoga para $\lambda(-x) = -(\lambda x)$.

\qed

\textit{Observação: Denotaremos $(-\lambda)x$ como $-\lambda x$ e $x+(-y)$ como $x-y$.}

\section{Subespaços vetoriais}

\begin{definition}

Sejam $\mathbb{K}$ um corpo, $E$ um $\mathbb{K}$-espaço vetorial e $F\subseteq E$ um subconjunto não vazio de $E$. Dizemos que $F$ é um \textbf{subespaço vetorial de $E$} se, e somente se, $F$ é um $\mathbb{K}$-espaço vetorial com as leis interna e externa de $E$.  

\end{definition}

\hypertarget{prop221}{}
\begin{proposition}

Sejam $\mathbb{K}$ um corpo, $E$ um $\mathbb{K}$-espaço vetorial e $F\subseteq E$ um subconjunto de $E$. Logo, $F$ é um subespaço vetorial de $E$ se, e somente se:

\begin{itemize}
  \item[(i)] $F \neq \emptyset$ 
  \item[(ii)] $x,y\in F \implies x+y\in F$
  \item[(iii)] $(x,\lambda)\in F\times\mathbb{K} \implies \lambda x \in F$
\end{itemize}

\end{proposition}

\textbf{Demonstração:}

Note que, se $F$ é subespaço vetorial de $E$, então evidentemente $0_E\in F$ e logo $F\neq\emptyset$ e para quaisquer $x,y\in F$ e $\lambda\in \mathbb{K}$ temos $x+y\in F$ e $\lambda x\in \mathbb{K}$. 

Suponha, agora, que $F\neq\emptyset$ e que para quaisquer $x,y\in F$ e $\lambda\in\mathbb{K}$ temos $x+y\in F$ e $\lambda x\in F$. Note que os axiomas referentes à lei externa de $E$ são trivialmente verificados para $F$, assim como a associatividade e comutatividade para a lei interna.

Resta mostrarmos que $0_E\in F$ e que para todo $x\in F$ existe $-x\in F$ tal que $x+(-x) = x-x = 0_E$.

Como $F\neq \emptyset$, então existe pelo menos um elemento $x\in F$. Além disso, como $\lambda x\in F$ para quaisquer $\lambda\in \mathbb{K}$, temos em particular que $0x = 0_E \in F$.

Analogamente, temos diretamente que $(-1)x = -x\in F$.

Portanto, concluímos que $F$ é um subespaço vetorial de $E$.

\qed

\textit{Observação:}

\begin{proposition}

Sejam $\mathbb{K}$ um corpo, $E$ um $\mathbb{K}$-espaço vetorial e $F\subseteq E$ um subconjunto de $E$. Logo, $F$ é um subespaço vetorial de $E$ se, e somente se:

\begin{itemize}
  \item[(i)] $F \neq \emptyset$ 
  \item[(ii)] $x,y\in F, \lambda,\mu\in \mathbb{K} \implies \lambda x+\mu y\in F$
\end{itemize}

\end{proposition}

\textbf{Demonstração:}

De forma análoga à demonstração da \hyperlink{prop221}{\textbf{Proposição 2.2.1}}, se $F$ for um subespaço vetorial de $E$, então é trivial que os itens $(i)$ e $(ii)$ são satisfeitos.

Suponha, agora, que os itens $(i)$ e $(ii)$ sejam verdadeiros. Analogamente à demonstração da \hyperlink{prop221}{\textbf{Proposição 2.2.1}}, basta mostrarmos que $0_E\in F$ e que para todo elemento $x\in F$ existe um inverso $-x\in F$.

Note que, para qualquer $x\in F$, temos:

\[
0_E = 0x + 0x\in F
\]

Logo, $0_E\in F$. Ainda, perceba que podemos reescrever $0_E$ como $0x$ para qualquer $x\in F$. Portanto, para um dado $x\in F$, temos ainda que:

\[
-x = (-1)x = 0_E + (-1)x = 0x + (-1)x \in F
\]

Logo, $F$ é subespaço vetorial de $E$.

\qed

\begin{proposition}
Sejam $\mathbb{K}$ um corpo, $E$ um $\mathbb{K}$-espaço vetorial e $\{F_i\}_{i\in I}$ uma família de subespaços vetoriais de $E$ indexados por $I$. Então:

\begin{itemize}
  \item[(i)] $\bigcap_{i\in I} F_i$ é subespaço vetorial de $E$.  
  \item[(ii)] $\bigcup_{i\in I} F_i$ é subespaço vetorial de $E$ se, e somente se, $\forall i,j\in I,\; \exists k\in I$ tal que $F_i\subseteq F_k$ e $F_j\subseteq F_k$ 
  \item[(iii)] $E\setminus F_i$ não é subespaço vetorial de $E$ para qualquer $i\in I$.
\end{itemize}

\end{proposition}

\textbf{Demonstração:}

Mostremos, primeiramente, o item $(i)$. Para isso, basta mostrarmos os itens da \hyperlink{prop221}{\textbf{Proposição 2.2.1}}.

Note que $0_E\in \bigcap_{i\in I} F_i$, trivialmente, uma vez que $F_i$ é subespaço vetorial de $E$ para todo $i\in I$.

Tomemos $x,y\in \bigcap_{i\in I} F_i$. Então, para todo $i\in I$, temos que $x,y\in F_i$ e, como $F_i$ é subespaço vetorial de $E$ para todo $i\in I$, deduzimos que $x+y\in F_i$, o que implica que $x+y\in \bigcap_{i\in I} F_i$.

A prova que $\lambda x \in \bigcap_{i\in I} F_i,\; \forall (\lambda, x)\in \mathbb{K}\times \bigcap_{i\in I} F_i$ é análoga.

Mostremos, agora, o item $(ii)$. Suponha inicialmente que para todo $i,j\in I,\;\exists k\in I$ tal que $F_i\subseteq F_k$ e $F_j\subseteq F_k$.

Note novamente que, trivialmente, temos $0_E\in \bigcup_{i\in I} F_i$.

Tomemos $x,y\in \bigcup_{i\in I} F_i$.

Por fim, para o item $(iii)$, note que $E\setminus F_i$ jamais será subespaço vetorial de $E$ para qualquer $i\in F_i$, pois $0_E \notin E\setminus F_i$. 

\qed

\end{document}
